Large language models are increasingly used to sort and rank data, yet we lack a systematic understanding of which properties they can reliably order by.
We introduce \benchname, a benchmark of 22 hand-curated ordering tasks spanning five property categories---syntactic, well-known factual, general knowledge, specific factual, and temporal---and evaluate two models (\gptlarge and \gptsmall) on one-pass listwise ordering.
We find that LLMs order items nearly perfectly ($\tau = 1.0$) when the target property corresponds to a well-known sequence (alphabetical order, planet distances, the periodic table) or a famous chronology (U.S.\ presidents, historical inventions), but struggle with continuous numeric properties that require precise recall of independent facts, such as city elevations ($\tau = 0.81$) and river lengths ($\tau = 0.88$).
\gptlarge significantly outperforms \gptsmall (mean $\tau = 0.948$ vs.\ $0.867$, Wilcoxon $p = 0.008$), yet both models agree on which properties are hardest (cross-model $\rho = 0.61$, $p = 0.003$).
These results reveal a property-accessibility hierarchy in parametric memory: algorithmic operations and memorized sequences are easiest, approximate magnitudes are intermediate, and precise numeric facts are hardest.
Our findings inform when to trust LLM-based data pipelines and highlight the boundary between categorical knowledge and precise quantitative recall.
